\chapter{Itô Calculus}\label{ch:qm:ito-calculus}

A junior researcher backtests a trend-following rule on ten years of daily closes: hold a
position proportional to the distance of the close above its 20-day moving average. The
Sharpe ratio is 5.4. The rule reads the closing price and the backtest credits it with the
return into that same close; shifted by one day, so that a position decided at a close
earns the next day's return, the Sharpe ratio is $-0.04$. The ``prices'' were a pure random
walk, with nothing to find. The first backtest was not an \omterm{def:qm:ito-calculus:integral}{Itô integral}: its integrand looked
at the increment it was multiplied by, and collected the covariation of the position with the
price, a quantity that exists on every random path. This chapter builds the \omterm{def:qm:ito-calculus:integral}{Itô integral},
whose integrand is fixed before the increment arrives, proves Itô's formula, the chain rule
with a second-order term that the rest of the series uses on every page, and ends with the
two structural theorems (Lévy's characterisation and \omterm{def:qm:probability-at-speed:martingale}{martingale} representation) that
underlie the pricing of One Quant Book~5.

\section{The Itô integral}

A trading strategy holds $H_t$ units of an asset whose price follows a \omterm{def:qm:brownian-motion:bm}{Brownian motion}; its
gains over $[0, T]$ should be $\int_0^T H_t\,dW_t$. The paths of $W$ have infinite total variation
(\cref{thm:qm:brownian-motion:qv}), so the integral cannot be defined path by path as a
Stieltjes integral, and the choice of evaluation point in each interval matters in the limit.
The \omterm{def:qm:ito-calculus:integral}{Itô integral} evaluates $H$ at the left point, which is the only choice a trader can make.

\begin{definition}[Simple process, Itô integral]\label{def:qm:ito-calculus:integral}
A \emph{simple process}\index{simple process} is $H_t = \sum_{k=0}^{n-1}h_k\mathbf 1_{(t_k,
t_{k+1}]}(t)$ with $0 = t_0 < \dots < t_n = T$ and each $h_k$ bounded and $\mathcal
F_{t_k}$-measurable. Its \emph{Itô integral}\index{Itô integral} is $\int_0^t H_s\,dW_s = \sum_k
h_k(W_{t_{k+1}\wedge t} - W_{t_k\wedge t})$. For an adapted $H$ with $\E\int_0^T H_s^2\,ds <
\infty$, the Itô integral is the $L^2$ limit of the integrals of simple processes $H^{(n)}$ with
$\E\int_0^T(H^{(n)}_s - H_s)^2\,ds \to 0$.
\end{definition}

\begin{theorem}[Itô isometry]\label{thm:qm:ito-calculus:isometry}
For adapted $H$ with $\E\int_0^T H_s^2\,ds < \infty$, $I_t = \int_0^t H_s\,dW_s$ is a continuous
square-integrable \omterm{def:qm:probability-at-speed:martingale}{martingale} with $I_0 = 0$ and
\[
\E\Bigl[\Bigl(\int_0^T H_s\,dW_s\Bigr)^2\Bigr] = \E\int_0^T H_s^2\,ds, \qquad [I]_t = \int_0^t H_s^2\,ds.
\]
\end{theorem}

\begin{proof}
For a \omterm{def:qm:ito-calculus:integral}{simple process}, with $\Delta_k = W_{t_{k+1}} - W_{t_k}$: the cross terms $\E[h_jh_k\Delta_j\Delta_k]$,
$j < k$, vanish because $h_jh_k\Delta_j$ is $\mathcal F_{t_k}$-measurable and $\E_{t_k}[\Delta_k] = 0$;
the square terms give $\E[h_k^2]\E[\Delta_k^2] = \E[h_k^2](t_{k+1} - t_k)$. The same computation on
$[s, t]$ gives the \omterm{def:qm:probability-at-speed:martingale}{martingale} property. The integral is therefore an isometry from simple
processes into $L^2$, extends to their closure, and the \omterm{def:qm:probability-at-speed:martingale}{martingale} property passes to $L^2$
limits (Doob's inequality, \cref{prop:qm:probability-at-speed:doob}, gives a continuous
version).
\end{proof}

\begin{example}[The integral of $W$ against itself]\label{ex:qm:ito-calculus:wdw}
On a partition of $[0, T]$, $\sum_k W_{t_k}\Delta_k = \tfrac12\sum_k(W_{t_{k+1}}^2 - W_{t_k}^2) -
\tfrac12\sum_k\Delta_k^2$. The first sum telescopes to $\tfrac12 W_T^2$ and the second tends to
$\tfrac12 T$ by \cref{thm:qm:brownian-motion:qv}:
\[
\int_0^T W_t\,dW_t = \tfrac12\bigl(W_T^2 - T\bigr).
\]
The right-point sums tend to $\tfrac12(W_T^2 + T)$, and the midpoint sums to $\tfrac12W_T^2$, the
ordinary calculus answer (\cref{fig:qm:ito-calculus:sums}).
\end{example}

\begin{omfigure}
\begin{tikzpicture}
\begin{axis}[width=11cm,height=5cm,xmode=log,log basis x=2,
  xlabel={points in $[0,1]$},ylabel={sum},
  tick label style={font=\small},label style={font=\small},
  xmin=3,xmax=90000,ymin=0.5,ymax=3.6,grid=major,grid style={gray!25},
  legend columns=3,legend style={font=\footnotesize,at={(0.5,-0.3)},anchor=north,draw=none,fill=none,/tikz/every even column/.append style={column sep=8pt}},
  table/col sep=comma]
\addplot[omThm,thick,mark=*,mark size=1.2pt] table[x=n,y=right]{figdata/methods/03-ito-calculus/sums.csv};
\addplot[omProp,thick,mark=square*,mark size=1.2pt] table[x=n,y=mid]{figdata/methods/03-ito-calculus/sums.csv};
\addplot[omDef,thick,mark=triangle*,mark size=1.6pt] table[x=n,y=left]{figdata/methods/03-ito-calculus/sums.csv};
\addplot[gray,dashed,forget plot] table[x=n,y=ito]{figdata/methods/03-ito-calculus/sums.csv};
\addplot[gray,dashed,forget plot] table[x=n,y=rlim]{figdata/methods/03-ito-calculus/sums.csv};
\legend{right point,midpoint (Stratonovich),left point (Itô)}
\end{axis}
\end{tikzpicture}

\omcaption{Sums approximating $\int_0^1 W\,dW$ on one path ($W_1 = 2.04$) sampled at $2^k$ points.
The left-point sums converge to $\tfrac12(W_1^2 - 1) = 1.58$, the right-point sums to $2.58$
(dashed lines); the midpoint sums equal $\tfrac12W_1^2 = 2.08$ exactly, whatever the grid. Data:
the chapter's tutorial, seeded.}\label{fig:qm:ito-calculus:sums}
\end{omfigure}

\begin{definition}[Stratonovich integral]\label{def:qm:ito-calculus:strat}
The \emph{Stratonovich integral}\index{Stratonovich integral} $\int_0^T X_t\circ dW_t$ is the limit of
the sums $\sum_k\tfrac12(X_{t_k} + X_{t_{k+1}})\Delta_k$. For a continuous semimartingale $X$ it
equals $\int_0^T X_t\,dW_t + \tfrac12[X, W]_T$.
\end{definition}

The \omterm{def:qm:ito-calculus:strat}{Stratonovich integral} obeys the ordinary chain rule, which is why physicists modelling a
smooth noise by a white one use it. It averages the integrand over the interval, which a trade
cannot do: the position is set before the price move. Finance uses Itô.

\section{The integral in a backtest}

A backtest on a grid $t_0 < t_1 < \dots$ with positions $\theta_{t_i}$ and prices $S_{t_i}$ computes
gains $\sum_i\theta_{t_i}(S_{t_{i+1}} - S_{t_i})$: the left-point sum, an \omterm{def:qm:ito-calculus:integral}{Itô integral}, when
$\theta_{t_i}$ uses only information available at $t_i$. A backtest that credits the position
computed at $t_{i+1}$ with the move into $t_{i+1}$ computes the right-point sum instead, and
\[
\sum_i\theta_{t_{i+1}}\Delta S_i - \sum_i\theta_{t_i}\Delta S_i = \sum_i\Delta\theta_i\,\Delta S_i
\longrightarrow [\theta, S]_T:
\]
the look-ahead profit is the covariation of the position with the price, which is not zero
on a random walk whenever the position responds to the latest price. The rule of the hook,
$\theta_t = (S_t - \bar S^{(20)}_t)/c$ with $\bar S^{(20)}$ the 20-day moving average and $c$ chosen
so that $\theta$ has unit variance, puts weight $(L - 1)/L = 0.95$ on the last increment:
\cref{fig:qm:ito-calculus:timing} shows the timing and \cref{fig:qm:ito-calculus:backtest} the
result.

\begin{omfigure}
\begin{tikzpicture}[x=3.4cm,y=1cm,font=\small]
\draw[->,thick,gray] (0,0) -- (3.3,0) node[right,black] {time};
\foreach \x/\l in {0/{close $t-1$},1/{close $t$},2/{close $t+1$}} {\draw[thick] (\x,-0.12) -- (\x,0.12); \node[below=3pt] at (\x,-0.1) {\l};}
\draw[omDef,very thick,-{Stealth}] (1,0.35) to[bend left=25] node[above,align=center] {Itô: $\theta_t$ earns\\$S_{t+1} - S_t$} (2,0.35);
\draw[omThm,very thick,-{Stealth},dashed] (1,-0.75) to[bend left=25] node[below=2pt,align=center] {look-ahead: $\theta_t$ earns\\$S_t - S_{t-1}$} (0,-0.75);
\node[draw=omDef,rounded corners=2pt,fill=omDef!6,inner sep=3pt,font=\footnotesize,align=center] (box) at (0.35,1.55) {position $\theta_t$ computed\\from the close $S_t$};
\draw[gray,->] (box.east) -| (1,0.18);
\end{tikzpicture}

\omcaption{Which price move a position may earn. A position computed from the close at $t$
earns the move from $t$ to $t + 1$ (solid): the left-point sum. Crediting it with the move into
$t$ (dashed) multiplies the position by an increment it has already seen.}\label{fig:qm:ito-calculus:timing}
\end{omfigure}

\begin{proposition}[The look-ahead Sharpe ratio on a random walk]\label{prop:qm:ito-calculus:lookahead}
Let $S$ be a random walk with daily increments $\sigma Z_t$, $Z_t$ independent standard normal,
and $\theta_t$ a unit-variance linear function of past increments with correlation $a$ with the
latest one. Then the same-bar gains $\theta_t\,\Delta S_t$ have mean $a\sigma$ and standard
deviation $\sigma\sqrt{1 + a^2}$, a daily Sharpe ratio of $a/\sqrt{1 + a^2}$, while the next-bar
gains have mean zero. For the moving-average rule, $a = \frac{(L-1)/L}{(\sum_{m=1}^{L-1}(m/L)^2)^{1/2}}$.
\end{proposition}

\begin{proof}
Write $\theta_t = aZ_t + \sqrt{1 - a^2}\,Y$ with $Y$ standard normal and independent of $Z_t$. Then
$\E[\theta_tZ_t] = a$ and $\E[\theta_t^2Z_t^2] = a^2\E[Z^4] + (1 - a^2) = 1 + 2a^2$, so the variance is
$1 + a^2$ (in units of $\sigma^2$). For the rule, $S_t - \bar S_t = \sum_{i=0}^{L-2}\frac{L-1-i}{L}\Delta
S_{t-i}$, whose variance is $\sigma^2\sum_{m=1}^{L-1}(m/L)^2$ and whose coefficient on $\Delta S_t$ is
$(L-1)/L$. Next-bar gains $\theta_{t-1}\Delta S_t$ have zero mean because $\theta_{t-1}$ is independent
of $\Delta S_t$.
\end{proof}

For $L = 20$, $a = 0.382$ and the annualised same-bar Sharpe ratio is $\sqrt{252}\times 0.382/
\sqrt{1.146} = 5.67$; over 400 simulated ten-year random walks it averages 5.69, and the honest
backtest 0.00 with a standard deviation of 0.32, the noise one expects from ten years of a
strategy with no edge.

\begin{omfigure}
\begin{tikzpicture}
\begin{axis}[width=11cm,height=5cm,
  xlabel={years},ylabel={cumulative P\&L (\% of unit notional)},
  tick label style={font=\small},label style={font=\small},
  xmin=0,xmax=10,ymin=-100,ymax=1050,grid=major,grid style={gray!25},
  legend columns=3,legend style={font=\footnotesize,at={(0.5,-0.3)},anchor=north,draw=none,fill=none,/tikz/every even column/.append style={column sep=8pt}},
  table/col sep=comma]
\addplot[omThm,very thick] table[x=year,y=cheat]{figdata/methods/03-ito-calculus/backtest.csv};
\addplot[omProp,thick,dashed] table[x=year,y=cov]{figdata/methods/03-ito-calculus/backtest.csv};
\addplot[omDef,very thick] table[x=year,y=honest]{figdata/methods/03-ito-calculus/backtest.csv};
\legend{same bar (look-ahead),covariation $[\theta{,}S]$,next bar (Itô)}
\end{axis}
\end{tikzpicture}

\omcaption{The 20-day trend rule on one ten-year random walk with 1\% daily volatility. Credited
with the same bar it earns 958\% of notional (Sharpe ratio 5.4), almost exactly the covariation of
position and price (965\%); credited with the next bar it earns $-7\%$ (Sharpe ratio $-0.04$).
Data: the chapter's tutorial, seeded.}\label{fig:qm:ito-calculus:backtest}
\end{omfigure}

\section{Itô's formula in one and several dimensions}

\begin{definition}[Itô process, quadratic covariation]\label{def:qm:ito-calculus:process}
An \emph{Itô process}\index{Itô process} is $X_t = X_0 + \int_0^t\mu_s\,ds + \int_0^t\sigma_s\,dW_s$, written
$dX_t = \mu_t\,dt + \sigma_t\,dW_t$, with adapted $\mu$, $\sigma$ such that $\int_0^T|\mu_s|\,ds$ and
$\int_0^T\sigma_s^2\,ds$ are finite almost surely. The \emph{quadratic covariation}\index{quadratic
covariation} of two continuous processes is the limit in probability $[X, Y]_t =
\lim\sum_k(X_{t_{k+1}} - X_{t_k})(Y_{t_{k+1}} - Y_{t_k})$; for Itô processes driven by $W^1, W^2$
with $d\langle W^1, W^2\rangle_t = \rho\,dt$, $d[X, Y]_t = \sigma^X_t\sigma^Y_t\rho\,dt$.
\end{definition}

The bookkeeping rules are $dt\,dt = 0$, $dt\,dW = 0$, $dW^i\,dW^j = \rho_{ij}\,dt$: only products of
two Brownian increments survive at first order in $dt$.

\begin{theorem}[Itô's formula]\label{thm:qm:ito-calculus:ito}
If $X$ is an \omterm{def:qm:ito-calculus:process}{Itô process} and $f \in C^{1,2}([0,T]\times\R)$, then
\[
df(t, X_t) = \partial_tf\,dt + \partial_xf\,dX_t + \tfrac12\partial_{xx}f\,d[X]_t
= \bigl(\partial_tf + \mu_t\partial_xf + \tfrac12\sigma_t^2\partial_{xx}f\bigr)dt + \sigma_t\partial_xf\,dW_t.
\]
For $X = (X^1, \dots, X^d)$ and $f \in C^{1,2}$, $df = \partial_tf\,dt + \sum_i\partial_if\,dX^i +
\tfrac12\sum_{i,j}\partial_{ij}f\,d[X^i, X^j]$; in particular $d(XY) = X\,dY + Y\,dX + d[X, Y]$.
\end{theorem}

\begin{proof}[Partial proof]
For $f(x)$ with bounded derivatives and $\mu, \sigma$ simple, Taylor's formula on a partition gives
$f(X_T) - f(X_0) = \sum f'(X_{t_k})\Delta X_k + \tfrac12\sum f^{\prime\prime}(\xi_k)(\Delta X_k)^2$. The first sum
converges to $\int f'(X)\,dX$. In the second, $f^{\prime\prime}(\xi_k)$ may be replaced by $f^{\prime\prime}(X_{t_k})$ at a
cost bounded by the modulus of continuity of $f^{\prime\prime}$ times $\sum(\Delta X_k)^2$, which tends to zero,
and $\sum f^{\prime\prime}(X_{t_k})((\Delta X_k)^2 - \sigma_{t_k}^2\Delta t_k)$ tends to zero in $L^2$ as in
\cref{thm:qm:brownian-motion:qv}. General $f$, $\mu$, $\sigma$ follow by localisation and
approximation (Karatzas and Shreve, 1991, §3.3).
\end{proof}

\begin{example}[The logarithm of a geometric Brownian motion]\label{ex:qm:ito-calculus:log}
For $dS = \mu S\,dt + \sigma S\,dW$, with $f = \ln$: $d\ln S = dS/S - \tfrac12 d[S]/S^2 = (\mu -
\tfrac12\sigma^2)dt + \sigma\,dW$. The expected return is $\mu$ but the expected log return, the rate
at which the median grows, is $\mu - \tfrac12\sigma^2$. With $\mu = 10\%$ and $\sigma = 40\%$, the
mean of $S_{10}/S_0$ is $e^{1} = 2.72$ and its median $e^{0.2} = 1.22$
(\cref{fig:qm:ito-calculus:gbm}); the volatility decay of leveraged funds (One Quant Book~1,
chapter~14) is the same correction applied twice.
\end{example}

\begin{omfigure}
\begin{tikzpicture}
\begin{axis}[width=11cm,height=5cm,ymode=log,
  xlabel={years},ylabel={$S_t/S_0$},
  tick label style={font=\small},label style={font=\small},
  xmin=0,xmax=10,ymin=0.2,ymax=8,grid=major,grid style={gray!25},
  ytick={0.25,0.5,1,2,4,8},log ticks with fixed point,
  legend columns=3,legend style={font=\footnotesize,at={(0.5,-0.3)},anchor=north,draw=none,fill=none,/tikz/every even column/.append style={column sep=8pt}},
  table/col sep=comma]
\addplot[name path=hi,gray!40,forget plot] table[x=t,y=p90]{figdata/methods/03-ito-calculus/gbm.csv};
\addplot[name path=lo,gray!40,forget plot] table[x=t,y=p10]{figdata/methods/03-ito-calculus/gbm.csv};
\addplot[gray!15,forget plot] fill between[of=hi and lo];
\addplot[omThm,very thick] table[x=t,y=mean_th]{figdata/methods/03-ito-calculus/gbm.csv};
\addplot[omDef,very thick] table[x=t,y=median_th]{figdata/methods/03-ito-calculus/gbm.csv};
\addplot[black,only marks,mark=*,mark size=1.3pt] table[x=t,y=median]{figdata/methods/03-ito-calculus/gbm.csv};
\legend{mean $e^{\mu t}$,median $e^{(\mu - \sigma^2/2)t}$,simulated median}
\end{axis}
\end{tikzpicture}

\omcaption{Geometric \omterm{def:qm:brownian-motion:bm}{Brownian motion} with $\mu = 10\%$ and $\sigma = 40\%$: the mean grows at 10\% a
year, the median at 2\%, and after ten years the 10th percentile (shaded band: 10th to 90th) is
0.24. Data: 100\,000 paths simulated exactly, seeded.}\label{fig:qm:ito-calculus:gbm}
\end{omfigure}

\section{Local martingales and the stochastic exponential}

An \omterm{def:qm:ito-calculus:integral}{Itô integral} $\int H\,dW$ with only $\int_0^TH_s^2\,ds < \infty$ almost surely need not be
integrable, let alone a \omterm{def:qm:probability-at-speed:martingale}{martingale}; it is one up to a sequence of \omterm{def:qm:probability-at-speed:stopping}{stopping times}.

\begin{definition}[Local martingale]\label{def:qm:ito-calculus:local}
An \omterm{def:qm:probability-at-speed:filtration}{adapted process} $M$ is a \emph{local martingale}\index{local martingale} if there are \omterm{def:qm:probability-at-speed:stopping}{stopping
times} $\tau_n \uparrow \infty$ such that each stopped process $M_{t\wedge\tau_n}$ is a \omterm{def:qm:probability-at-speed:martingale}{martingale}.
\end{definition}

\begin{proposition}[Positive local martingales are supermartingales]\label{prop:qm:ito-calculus:super}
A \omterm{def:qm:ito-calculus:local}{local martingale} bounded below, in particular a positive one, is a \omterm{def:qm:probability-at-speed:martingale}{supermartingale}. A
\omterm{def:qm:ito-calculus:local}{local martingale} $M$ with $\E[\sup_{s\le t}|M_s|] < \infty$ for every $t$ is a \omterm{def:qm:probability-at-speed:martingale}{martingale}.
\end{proposition}

\begin{proof}
If $M \ge 0$, Fatou's lemma gives $\E_s[M_t] = \E_s[\lim_n M_{t\wedge\tau_n}] \le \liminf_n
M_{s\wedge\tau_n} = M_s$. Under the domination, dominated convergence replaces the inequality by an
equality.
\end{proof}

The inequality can be strict: $1/|B_t|$ for a three-dimensional \omterm{def:qm:brownian-motion:bm}{Brownian motion} started away from
the origin is a positive \omterm{def:qm:ito-calculus:local}{local martingale} whose expectation decreases, a \emph{strict} \omterm{def:qm:ito-calculus:local}{local
martingale}. Such processes model price bubbles and are why chapter~5 checks that a candidate
\omterm{def:qm:probability-at-speed:change}{density process} is a true \omterm{def:qm:probability-at-speed:martingale}{martingale} before using it to change the measure.

\begin{definition}[Stochastic exponential]\label{def:qm:ito-calculus:exp}
The \emph{stochastic exponential}\index{stochastic exponential} of a continuous semimartingale $X$
with $X_0 = 0$ is $\mathcal E(X)_t = \exp(X_t - \tfrac12[X]_t)$, the unique solution of $dZ = Z\,dX$,
$Z_0 = 1$.
\end{definition}

By Itô's formula, $d\exp(X - \tfrac12[X]) = \exp(\cdot)(dX - \tfrac12 d[X] + \tfrac12 d[X]) = Z\,dX$. With
$X = \int\gamma\,dW$ it is the positive \omterm{def:qm:ito-calculus:local}{local martingale} $\exp(\int\gamma\,dW - \tfrac12\int\gamma^2dt)$,
the \omterm{def:qm:probability-at-speed:change}{density process} of every \omterm{def:qm:probability-at-speed:change}{change of measure} in chapter~5.

\section{The representation theorem}

\begin{theorem}[Lévy's characterisation]\label{thm:qm:ito-calculus:levy}
A continuous \omterm{def:qm:ito-calculus:local}{local martingale} $M$ with $M_0 = 0$ and $[M]_t = t$ is a \omterm{def:qm:brownian-motion:bm}{Brownian motion}.
\end{theorem}

\admitted

\begin{theorem}[Martingale representation]\label{thm:qm:ito-calculus:representation}
Let $\mathbb F$ be the \omterm{def:qm:probability-at-speed:filtration}{filtration} generated by a \omterm{def:qm:brownian-motion:bm}{Brownian motion} $W$. Every square-integrable
$\mathbb F$-martingale $M$ has the form $M_t = M_0 + \int_0^tH_s\,dW_s$ for a unique adapted $H$ with
$\E\int_0^TH_s^2\,ds < \infty$.
\end{theorem}

\admitted

Both are proved in Karatzas and Shreve (1991, §3.3 and §3.4). The representation theorem is the
mathematics of replication: if the only randomness is $W$, every payoff's \omterm{def:qm:probability-at-speed:condexp}{conditional expectation}
is a stochastic integral against $W$, and $H$ is the position that reproduces it. One Quant
Book~5, chapter~1, builds on it the notions of a replicating portfolio and a complete market; it
fails as soon as there are jumps (chapter~6) or more sources of risk than traded assets.

\section{Tutorial: Itô sums and an honest backtest}\label{tut:qm:ito-calculus:backtest}

\begin{tutorial}
\textbf{Goal.} See the left point of the \omterm{def:qm:ito-calculus:integral}{Itô integral} in numbers, check Itô's formula on one
path, and measure the look-ahead profit of a same-bar backtest. \textbf{End state:}
\cref{fig:qm:ito-calculus:sums,fig:qm:ito-calculus:backtest,fig:qm:ito-calculus:gbm} and the Sharpe
ratios 5.4 and $-0.04$.
\begin{tutsteps}
\item \textbf{The integrals.} The running project's gains are left-point sums; the same-bar
version and the covariation that separates them are one line each.
\omcode{code/firm/stochint/firm_stochint.py}{24}{42}{Adapted gains, same-bar gains and the covariation between them.}\label{lst:qm:ito-calculus:gains}
\item \textbf{The rule}: the distance of the close from its moving average, scaled to unit
variance under a random walk, and the look-ahead mean of \cref{prop:qm:ito-calculus:lookahead}.
\omcode{code/methods/03-ito-calculus/python/qm_ito.py}{33}{47}{The trend signal and the expected look-ahead gain.}\label{lst:qm:ito-calculus:signal}
\item \textbf{Run} \texttt{sums\_table()}, \texttt{ito\_check()} (on a GBM path with $2^{14}$ steps,
$\ln S_1 = 0.70040$ against $0.70039$ from Itô's formula and $0.78179$ without the $-\tfrac12\,d[S]/S^2$
term), \texttt{backtest()} and \texttt{fig\_ito.py}.
\end{tutsteps}
\textbf{What to change next.} Replace the trend rule by a mean-reversion rule and predict the sign
of the look-ahead bias; add a one-tick bid--ask bounce to the prices and see what the honest
backtest of a mean-reversion rule then reports.
\end{tutorial}

\section{Build: discrete stochastic integrals}\label{bld:qm:ito-calculus:stochint}

\begin{build}
\bfield{Purpose} Every backtest of the miniature firm computes its P\&L through these functions,
so that a position can earn only the price move after it is decided, and a look-ahead leaves a
measurable trace.
\bfield{Interface} \texttt{gains(theta, s)}; \texttt{same\_bar\_gains(theta, s)};
\texttt{covariation(x, y)}; \texttt{lookahead\_test(theta, s)} returning the same-bar excess, the
covariation it equals, a $t$-statistic and the identity error; \texttt{riemann\_sums(h, x,
point)}.
\bfield{Rules} Positions and prices are aligned on the same timestamps and the function, not
the caller, shifts them; the identity same-bar $-$ adapted $=$ covariation holds to rounding.
\bfield{Acceptance tests} \texttt{code/firm/stochint/tests/}: on random walks the adapted gains of
any rule have mean zero; the same-bar gains of the trend rule have the mean of
\cref{prop:qm:ito-calculus:lookahead}; the identity holds to $10^{-12}$; the three Riemann sums of
$\int W\,dW$ converge to their limits.
\bfield{Stretch} Timestamps with latencies (a position becomes live some microseconds after its
signal); transaction costs as a function of $|\Delta\theta|$.
\end{build}

\begin{omsources}
\item K.~Itô, ``Stochastic integral'', \emph{Proceedings of the Imperial Academy, Tokyo} 20, 1944.
\item K.~Itô, ``On a formula concerning stochastic differentials'', \emph{Nagoya Mathematical
Journal} 3, 1951.
\item R.~L.~Stratonovich, ``A new representation for stochastic integrals and equations'',
\emph{SIAM Journal on Control} 4, 1966.
\item I.~Karatzas and S.~E.~Shreve, \emph{Brownian Motion and Stochastic Calculus}, Springer,
2nd ed., 1991.
\end{omsources}

\section{Exercises}

\begin{exercise}[$\star$]\label{exo:qm:ito-calculus:1}
Compute $\E[(\int_0^1W_t\,dW_t)^2]$ with the isometry and check it with \cref{ex:qm:ito-calculus:wdw}.
\end{exercise}

\begin{exercise}[$\star$]\label{exo:qm:ito-calculus:2}
Write $d(W_t^3)$ with Itô's formula, and deduce $\E[W_t^3]$ and a \omterm{def:qm:probability-at-speed:martingale}{martingale} built from $W^3$.
\end{exercise}

\begin{exercise}[$\star$]\label{exo:qm:ito-calculus:3}
A stock follows a geometric \omterm{def:qm:brownian-motion:bm}{Brownian motion} with $\mu = 10\%$ and $\sigma = 40\%$. What are the mean
and the median of $S_{10}/S_0$, and at what volatility would the median not grow at all?
\end{exercise}

\begin{exercise}[$\star\star$]\label{exo:qm:ito-calculus:4}
Show that $e^{W_t - t/2}$ is a \omterm{def:qm:probability-at-speed:martingale}{martingale} and compute $\E[e^{2W_1}]$.
\end{exercise}

\begin{exercise}[$\star\star$]\label{exo:qm:ito-calculus:5}
Use the product rule on $tW_t$ to write $\int_0^1 W_t\,dt$ as a stochastic integral, and compute its
variance.
\end{exercise}

\begin{exercise}[$\star\star$]\label{exo:qm:ito-calculus:6}
$X = W^1$ and $Y = 2W^1 + W^2$ with $W^1, W^2$ independent. Compute $[X, Y]_t$, $[Y]_t$ and the
instantaneous correlation of $X$ and $Y$.
\end{exercise}

\begin{exercise}[$\star\star\star$]\label{exo:qm:ito-calculus:7}
\emph{Coding.} Repeat the backtest with a 60-day moving average: compute the look-ahead Sharpe
ratio from \cref{prop:qm:ito-calculus:lookahead} and its average over 100 simulated ten-year walks.
\end{exercise}

\begin{exercise}[$\star\star\star$]\label{exo:qm:ito-calculus:8}
\emph{Find the flaw.} ``Our rule decides from the closing price and trades in the closing
auction; the backtest credits it with the next day's return, so there is no look-ahead.''
\end{exercise}

\section{Problem: The Look-Ahead in the Backtest}

\begin{problem}[{Weekend problem --- a Sharpe ratio of 5.4 on a random walk}]\label{pb:qm:ito-calculus:1}
The log price is a random walk with 1\% daily volatility. The rule holds $\theta_t = (S_t -
\bar S^{(20)}_t)/c$, with $c$ making $\theta$ of unit variance; the flawed backtest credits $\theta_t$
with $S_t - S_{t-1}$, the honest one with $S_{t+1} - S_t$.

\medskip\noindent\textbf{Part I --- The rule.}
\begin{enumerate}
\item What is $c$?
\item What is the correlation $a$ of $\theta_t$ with the day's increment?
\item What are the mean and standard deviation of the same-bar gain per day?
\item What is the annualised same-bar Sharpe ratio?
\item What does the honest backtest earn on average, and why?
\end{enumerate}

\medskip\noindent\textbf{Part II --- The simulation.}
\begin{enumerate}[resume]
\item What Sharpe ratios does the seeded ten-year history give for the two backtests?
\item Over 400 histories, what are the mean and standard deviation of each?
\item Check the identity same-bar $-$ honest $=$ covariation on the seeded history.
\item What $t$-statistic does \texttt{lookahead\_test} report for the seeded history?
\item Why is the honest Sharpe ratio's standard deviation about $1/\sqrt{10}$?
\end{enumerate}

\medskip\noindent\textbf{Part III --- The Itô view.}
\begin{enumerate}[resume]
\item Write the same-bar P\&L as an \omterm{def:qm:ito-calculus:integral}{Itô integral} plus a covariation.
\item Why is the expectation of the Itô part zero?
\item What are the look-ahead Sharpe ratios for windows of 5, 10 and 60 days?
\item If the flawed backtest ran on five-minute bars, 78 a day, with the same rule in bars, what would its annualised Sharpe ratio be?
\item Would trading at the next day's open instead of the close remove the bias?
\end{enumerate}

\medskip\noindent\textbf{Part IV --- Judgement.}
\begin{enumerate}[resume]
\item How would you detect this error in someone else's backtest?
\item Is look-ahead bias always positive?
\item On real prices with a bid--ask bounce, what else can produce a covariation between position and price?
\item State the \emph{named result}: the look-ahead Sharpe ratio of the same-bar trend rule on a random walk.
\item In one sentence: what makes a backtest an \omterm{def:qm:ito-calculus:integral}{Itô integral}?
\end{enumerate}
\end{problem}

\section{Interview questions}

\begin{interviewq}[$\star$ \iqroles{researcher}]\label{iq:qm:ito-calculus:1}
Compute $\int_0^TW_t\,dW_t$. Why is it not $\tfrac12W_T^2$?
\end{interviewq}

\begin{interviewq}[$\star$ \iqroles{trader, researcher}]\label{iq:qm:ito-calculus:2}
A stock has an expected return of 10\% and a volatility of 40\%. What is its expected log return,
and what does a typical holder earn over ten years?
\end{interviewq}

\begin{interviewq}[$\star\star$ \iqroles{researcher}]\label{iq:qm:ito-calculus:3}
Is $W_t^2$ a \omterm{def:qm:probability-at-speed:martingale}{martingale}? Find the process you must subtract, and the variance of $\int_0^TW\,dW$.
\end{interviewq}

\begin{interviewq}[$\star\star$ \iqroles{developer, researcher}]\label{iq:qm:ito-calculus:4}
How would you design a backtesting engine that cannot commit look-ahead by construction?
\end{interviewq}

\begin{interviewq}[$\star\star$ \iqroles{researcher, bank}]\label{iq:qm:ito-calculus:5}
What is a \omterm{def:qm:ito-calculus:local}{local martingale} that is not a \omterm{def:qm:probability-at-speed:martingale}{martingale}, and why should a pricing quant care?
\end{interviewq}

\begin{interviewq}[$\star\star\star$ \iqroles{researcher, mle}]\label{iq:qm:ito-calculus:6}
Itô or Stratonovich: which does finance use, which does physics use, and why?
\end{interviewq}
